\section*{Hoofdstuk \ref{ch:b1:enzymes} --- Enzymen en biochemische katalyse}

\begin{solution}{exo:b1:enzymes:1}
Het verlaagt de \omterm{prop:b1:enzymes:whatitdoes}{activeringsenergie} (de hoogte van de \omterm{prop:b1:enzymes:whatitdoes}{overgangstoestand}), en
daarmee de snelheid, in beide richtingen; het laat $\Delta G$, de
evenwichtsconstante en de ligging van het evenwicht ongemoeid. Op het profiel
wordt de piek verlaagd, de twee uiteinden niet.
\end{solution}

\begin{solution}{exo:b1:enzymes:2}
$K_m$: de substraatconcentratie bij de helft van $V_{\max}$ (\unit{mol/L}).
$V_{\max}$: de snelheid bij verzadigend \omterm{def:b1:enzymes:enzyme}{substraat} (\unit{mol/s}, of \unit{\micro
mol/min}). $k_{\text{cat}}$: het aantal substraatmoleculen dat per seconde per
\omterm{def:b1:enzymes:enzyme}{enzym} bij verzadiging wordt omgezet (\unit{s^{-1}}). $k_{\text{cat}}/K_m$: de
efficiëntie, de tweede-ordesnelheidsconstante bij weinig \omterm{def:b1:enzymes:enzyme}{substraat}
(\unit{L.mol^{-1}.s^{-1}}).
\end{solution}

\begin{solution}{exo:b1:enzymes:3}
Zonder \omterm{def:b1:enzymes:inhibitors}{remmer}: snijpunt 1, dus $V_{\max} = 1$; snijpunt met de $x$-as $-0.5$, dus
$K_m = 2$. Competitief: $V_{\max} = 1$, snijpunt met de $x$-as $-0.25$, $K_m = 4$.
Niet-competitief: snijpunt 2, $V_{\max} = 0.5$; snijpunt met de $x$-as $-0.5$,
$K_m = 2$.
\end{solution}

\begin{solution}{exo:b1:enzymes:4}
Allosterische (feedback)regeling, milliseconden; \omterm{prop:b1:enzymes:control}{covalente modificatie}
(fosforylering), seconden tot minuten; proteolytische activering van een zymogeen,
eenmalig en op afroep; sturing van de hoeveelheid via genexpressie en afbraak,
minuten tot uren.
\end{solution}

\begin{solution}{exo:b1:enzymes:5}
$v_0 = 50[S]/(0.1 + [S])$: 8.3, 25, 45.5 en \qty{49.5}{\micro mol/min}.
$0.9 = [S]/(0.1 + [S])$ geeft $[S] = 0.9\,K_m/0.1 = \qty{0.9}{mmol/L}$.
\end{solution}

\begin{solution}{exo:b1:enzymes:6}
$e^{20\,000/2580} = e^{7.75} = 2300$. Voor $10^{10}$: $\Delta E_a =
RT\ln 10^{10} = 2.58\times 23.0 = \qty{59}{kJ/mol}$.
\end{solution}

\begin{solution}{exo:b1:enzymes:7}
$V_{\max} = \qty{0.12}{mol/L}$ per minuut in \qty{1}{mL}:
\qty{1.2e-4}{mol/min} $= \qty{2.0e-6}{mol/s}$ product; \omterm{def:b1:enzymes:enzyme}{enzym} \qty{2e-12}{mol}:
$k_{\text{cat}} = 2\times 10^{-6}/2\times 10^{-12} =
\qty{1e6}{s^{-1}}$. Efficiëntie $10^6/0.012 = \qty{8e7}{L.mol^{-1}.s^{-1}}$,
binnen een factor tien van de diffusiegrens: vrijwel elke ontmoeting met een
substraatmolecuul is vruchtbaar.
\end{solution}

\begin{solution}{exo:b1:enzymes:8}
$1 + 2/K_i = 2$: $K_i = \qty{2}{mmol/L}$. Voor \qty{90}{\%}:
$[S] = 9K_m^{\text{app}}$: \qty{9}{mmol/L} zonder, \qty{18}{mmol/L} met de \omterm{def:b1:enzymes:inhibitors}{remmer}.
\end{solution}

\begin{solution}{exo:b1:enzymes:9}
Het remmen van de eerste vastleggende stap legt de hele route in één keer stil en
verspilt geen tussenproducten, die zich anders zouden ophopen; het eindproduct is
het signaal dat ertoe doet --- de overvloed ervan is precies de informatie dat de
route niet langer nodig is --- terwijl het gehalte van een tussenproduct niets
over de vraag zegt.
\end{solution}

\begin{solution}{exo:b1:enzymes:10}
Een protease dat actief zou zijn in de \omterm{def:b1:cell-unit-of-life:cell}{cel} die het maakt, zou die \omterm{def:b1:cell-unit-of-life:cell}{cel} verteren;
als inactief zymogeen gemaakt is het onschadelijk tot enteropeptidase, een \omterm{def:b1:enzymes:enzyme}{enzym}
van de darmbekleding, trypsinogeen in de twaalfvingerige darm tot trypsine knipt,
waar de vertering gewenst is; trypsine activeert vervolgens de overige zymogenen
(een cascade). Omdat een spoortje trypsine de cascade te vroeg op gang kan
brengen, verpakt de alvleesklier voor de zekerheid een specifieke trypsineremmer
in haar korrels; faalt die regeling, dan volgt een alvleesklierontsteking.
\end{solution}

\begin{solution}{exo:b1:enzymes:11}
Zonder: $3^3/(3^3 + 3^3) = 0.50$. Met: $27/(166 + 27) = 0.14$: de snelheid daalt
tot \qty{28}{\%} van haar waarde. Bij een michaelis-mentenenzym met $K_m$ verhoogd
van 3 naar 5.5: van $3/6 = 0.50$ naar $3/8.5 = 0.35$, oftewel \qty{70}{\%}.
\omterm{def:b1:proteins:allostery}{Coöperativiteit} maakt dezelfde verschuiving van de curve rond het werkpunt meer
dan twee keer zo doeltreffend: een schakelaar in plaats van een dimmer.
\end{solution}

\begin{solution}{exo:b1:enzymes:12}
Een kale katalysator (een platina-oppervlak, een \omterm{def:b1:water-small-molecules:ph}{zuur}) versnelt zonder onderscheid
vele reacties; een \omterm{def:b1:enzymes:enzyme}{enzym} versnelt één reactie op één \omterm{def:b1:enzymes:enzyme}{substraat} --- de specificiteit
van zijn \omterm{def:b1:enzymes:enzyme}{actieve plaats} is al informatie over wat de \omterm{def:b1:cell-unit-of-life:cell}{cel} gedaan wil hebben. De
regeling voegt een tweede laag toe: allosterische plaatsen, fosforyleringsplaatsen
en zymogeenpeptiden vertellen het \omterm{def:b1:enzymes:enzyme}{enzym} wanneer en waar het moet werken, naar
gelang de toestand van de \omterm{def:b1:cell-unit-of-life:cell}{cel}. De kosten zijn een groot \omterm{def:b1:proteins:peptide}{eiwit} (honderden residuen
voor een plaats van een tiental), een tragere omzet dan bij de beste anorganische
katalysatoren, en de machinerie om het te maken, te wijzigen en af te breken; de
winst is een chemie die alleen loopt waar en wanneer zij nuttig is, en dat is wat
een stofwisseling is.
\end{solution}

\begin{solution}{pb:b1:enzymes:1}
\textbf{1.} $1/[S]$: 2, 1, 0.5, 0.2, 0.1, 0.05. $1/v_0$: 5.0, 3.0, 2.0, 1.4,
1.2, 1.1.
\textbf{2.} Elke stap in $1/[S]$ verandert $1/v_0$ evenredig: helling
$(5.0 - 1.1)/(2 - 0.05) = 2.0$; snijpunt $1.0$.
\textbf{3.} $V_{\max} = 1/1.0 = \qty{1.0}{\micro mol/min}$; $K_m =
\text{helling}\times V_{\max} = \qty{2.0}{mmol/L}$.
\textbf{4.} $1.0\times 5/(2 + 5) = 0.714$: zoals gemeten.
\textbf{5.} $\qty{10e-9}{mol/L}\times\qty{1e-3}{L} = \qty{1e-11}{mol}$;
$V_{\max} = \qty{1e-6}{mol/min} = \qty{1.67e-8}{mol/s}$;
$k_{\text{cat}} = 1.67\times 10^{-8}/10^{-11} = \qty{1670}{s^{-1}}$.
\textbf{6.} $1670/(2\times 10^{-3}) = \qty{8.3e5}{L.mol^{-1}.s^{-1}}$: duizend
keer onder de diffusiegrens; slechts één botsing op de duizend is vruchtbaar.
\textbf{7.} $0.95 = [S]/(2 + [S])$: $[S] = 19\,K_m = \qty{38}{mmol/L}$.
\textbf{8.} $1/v_0$: 9.0, 5.0, 3.0, 1.8, 1.4, 1.2. Helling $(9.0 -
1.2)/1.95 = 4.0$; snijpunt $1.0$.
\textbf{9.} $V_{\max}^{\text{app}} = \qty{1.0}{\micro mol/min}$ (onveranderd);
$K_m^{\text{app}} = 4.0\times 1.0 = \qty{4.0}{mmol/L}$.
\textbf{10.} $K_m$ verdubbeld, $V_{\max}$ onveranderd: competitief.
\textbf{11.} $4 = 2(1 + 1/K_i)$: $K_i = \qty{1.0}{mmol/L}$.
\textbf{12.} Met de \omterm{def:b1:enzymes:inhibitors}{remmer} geldt $v = V_{\max}[S]/(K_m(1 + [I]/K_i) +
[S])$; de snelheid halveren vraagt $K_m(1 + [I]/K_i) + [S] = 2(K_m +
[S])$, oftewel $[I] = K_i(K_m + [S])/K_m$: bij \qty{2}{mmol/L} is $[I] =
1\times 4/2 = \qty{2}{mmol/L}$; bij \qty{20}{mmol/L} is $[I] = 22/2 =
\qty{11}{mmol/L}$. Hoe meer \omterm{def:b1:enzymes:enzyme}{substraat}, hoe meer \omterm{def:b1:enzymes:inhibitors}{remmer} er nodig is.
\textbf{13.} Een molecuul dat op het \omterm{def:b1:enzymes:enzyme}{substraat} (of op zijn \omterm{prop:b1:enzymes:whatitdoes}{overgangstoestand})
lijkt --- een esteranaloog dat niet kan worden gehydrolyseerd --- en dat in de
\omterm{def:b1:enzymes:enzyme}{actieve plaats} bindt.
\textbf{14.} $1/v_0$: 10.0, 6.0, 4.0, 2.8, 2.4, 2.2; helling $4.0$, snijpunt
$2.0$: $V_{\max}^{\text{app}} = \qty{0.5}{\micro mol/min}$,
$K_m^{\text{app}} = 4.0\times 0.5 = \qty{2.0}{mmol/L}$.
\textbf{15.} $V_{\max}$ gehalveerd, $K_m$ onveranderd: niet-competitief;
$2 = 1 + 1/K_i$: $K_i = \qty{1.0}{mmol/L}$.
\textbf{16.} J bindt op een andere plaats dan de actieve, zowel op het vrije \omterm{def:b1:enzymes:enzyme}{enzym}
als op $ES$, met dezelfde affiniteit: \omterm{def:b1:enzymes:enzyme}{substraat} concurreert niet met hem, en bij
elke $[S]$ is de helft van de enzymmoleculen uitgeschakeld.
\textbf{17.} Met J verdubbelt de snelheid (de helft van twee keer zoveel \omterm{def:b1:enzymes:enzyme}{enzym} is
nog altijd actief); met I verdubbelt zij ook (de snelheid is bij elke $[S]$
evenredig met het \omterm{def:b1:enzymes:enzyme}{enzym}): het verdubbelen van het \omterm{def:b1:enzymes:enzyme}{enzym} onderscheidt de \omterm{def:b1:enzymes:inhibitors}{remmers}
nooit.
\textbf{18.} Snelheid bij \qty{25}{\celsius}: $V_{\max}[S]/(K_m + [S])$ met
$V_{\max} = k_{\text{cat}}[E]$: per liter $1670\times 10^{-8} =
\qty{1.67e-5}{mol/L/s}$; $\times 0.5/2.5 = \qty{3.3e-6}{mol/L/s}$; bij
\qty{37}{\celsius} $\times 2^{1.2} = 2.3$: \qty{7.7e-6}{mol/L/s}; in
\qty{1e-12}{L}: \qty{7.7e-18}{mol/s} $= \num{4.6e6}$ moleculen per seconde.
\textbf{19.} Er is geen verhoging nodig: de snelheid (\num{4.6e6}) overtreft de
vraag (\num{3e6}); moest zij toch stijgen, dan kon de \omterm{def:b1:cell-unit-of-life:cell}{cel} het \omterm{def:b1:enzymes:enzyme}{substraat} verhogen,
het \omterm{def:b1:enzymes:enzyme}{enzym} fosforyleren om zijn $K_m$ te verlagen, of meer \omterm{def:b1:enzymes:enzyme}{enzym} maken.
\textbf{20.} $K_m^{\text{app}} = 2(1 + 0.4/0.2) = \qty{6}{mmol/L}$: de snelheid
wordt $\times 0.5/6.5$ in plaats van $0.5/2.5$: \qty{38}{\%} van daarvoor,
\num{1.8e6} per seconde --- het product knijpt zijn eigen synthese af tot onder de
vraag, zodat het wordt verbruikt en zijn gehalte daalt, wat het \omterm{def:b1:enzymes:enzyme}{enzym} weer
vrijgeeft: een zichzelf bijstellende aanvoer.
\textbf{21.} Vóór: $0.5/2.5 = 0.20$ van $V_{\max}$; na: $0.5/0.9 =
0.56$: bijna drievoudig.
\textbf{22.} Vrijwel nul: bij pH 5 is de histidine van de \omterm{def:b1:enzymes:enzyme}{actieve plaats}
geprotoneerd en kan zij niet als base werken, en zijn de zure residuen van de
plaats geneutraliseerd; de ionisatietoestand van het \omterm{def:b1:enzymes:enzyme}{enzym}, en misschien ook zijn
vouwing, deugen niet. Het is voor het \omterm{def:b1:eukaryotic-cell:organelle}{cytosol} gebouwd, niet voor het \omterm{def:b1:eukaryotic-cell:endomembrane}{lysosoom}.
\textbf{23.} $K_m$ verandert weinig (het binden gebeurt grotendeels door de rest
van de plaats); $k_{\text{cat}}$ stort met ordes van grootte in, want de histidine
was de katalytische base die het water of de serine activeerde.
\textbf{24.} Dicht bij $K_m$ reageert de snelheid op veranderingen van het
\omterm{def:b1:enzymes:enzyme}{substraat} (de helling ligt rond $V_{\max}/2K_m$); ver erboven is het \omterm{def:b1:enzymes:enzyme}{enzym}
verzadigd en ongevoelig, en zou het overtollige \omterm{def:b1:enzymes:enzyme}{substraat} een osmotische en
chemische last zijn. Dicht bij $K_m$ werken maakt de route via de aanvoer
stuurbaar.
\textbf{25.} $K_m = \qty{2.0}{mmol/L}$, $V_{\max} = \qty{1.0}{\micro mol/min}$
(bij \qty{10}{nmol/L} \omterm{def:b1:enzymes:enzyme}{enzym}), $k_{\text{cat}} = \qty{1670}{s^{-1}}$; I
competitief, $K_i = \qty{1.0}{mmol/L}$; J niet-competitief, $K_i =
\qty{1.0}{mmol/L}$.
\end{solution}
